\chapter{Computational Complexity and Decidability Bounds}
\label{ch:computational_complexity}

In this chapter, we rigorously establish the foundational computational bounds governing our automated synthesis frameworks. The preceding chapters introduced the syntactic and semantic profiles of the Praxis engine, laying the groundwork for how agents communicate, store state, and parse instructions. We now turn our attention to the theoretical limits of computation, mapping the state-space exploration of agentic interactions to classical Turing machine semantics and formal decidability algebra. 

The advent of Large Language Models (LLMs) and agentic loops introduces a paradigm where computation is bounded not just by classical time and space limitations, but by the \textit{context window}—a strict physical limit on the number of tokens an agent can process simultaneously. Our primary objective in this chapter is to formalize these boundaries. We will prove the Halting bounds within bounded-context execution, demonstrating that isolated agent operations are strictly decidable. We will then expand our analysis to the complexity of the \textit{receipt verification} protocol under BLAKE3 hashing constraints, establishing the algebraic structures that guarantee determinism. Furthermore, we analyze the architectural mandate for linear scalability, the mathematical topology of the Refusal Manifold in closed vocabularies, and prove strict adherence to the symmetry of space complexity demanded by the core team discipline.

By embedding our system's physical constraints into formal automata theory, we bridge the gap between empirical software engineering—such as our Rust-based core discipline—and the foundational mathematics of computation.

\section{Formal Setup: Turing Machines and Bounded Contexts}

The standard Turing machine assumes an infinite tape, rendering the general Halting Problem undecidable. However, real-world agentic systems operate under strict token limits. To model this, we introduce the Context-Bounded Turing Machine (CBTM), which structurally enforces a hard limit on the active memory space, analogous to the finite context window of an LLM.

\begin{definition}[Context-Bounded Turing Machine (CBTM)]
\label{def:cbtm}
A Context-Bounded Turing Machine is a 7-tuple $M = (Q, \Sigma, \Gamma, \delta, q_0, q_{accept}, q_{reject}, \kappa)$ where:
\begin{itemize}
    \item $Q$ is a finite, non-empty set of internal states, reflecting the operational phases of the agent's parsing and reasoning engine.
    \item $\Sigma$ is the finite input alphabet, strictly excluding the blank symbol $\sqcup$. In our architecture, $\Sigma$ corresponds to the set of valid vocabulary tokens.
    \item $\Gamma$ is the finite tape alphabet, representing both input tokens and intermediate scratchpad embeddings, where $\Sigma \subset \Gamma$ and $\sqcup \in \Gamma$.
    \item $\delta: Q \times \Gamma \rightarrow Q \times \Gamma \times \{L, R\}$ is the deterministic transition function governing state changes, symbol writing, and tape head movement.
    \item $q_0 \in Q$ is the uniquely designated start state.
    \item $q_{accept}, q_{reject} \in Q$ are the distinct halting states ($q_{accept} \neq q_{reject}$).
    \item $\kappa \in \mathbb{N}$ represents the strict upper bound on the number of contiguous tape cells that can be actively modified or traversed—formally, the context bound.
\end{itemize}
\end{definition}

The introduction of $\kappa$ models the physical constraints of contemporary agentic windows (e.g., a $128,000$ token limit). The computation of $M$ on input string $w \in \Sigma^*$ proceeds conventionally, with one critical structural enforcement: if the transition function $\delta$ commands the tape head to transition to the $(\kappa+1)$-th distinct cell relative to the initial position, the computation immediately aborts. The machine undergoes a forced state transition to $q_{reject}$, emitting a context-overflow refusal. 

To analyze the complexity classes applicable to CBTMs, we first must quantify their state space.

\begin{lemma}[Finiteness of Configuration Space]
\label{lem:finite_config}
For any Context-Bounded Turing Machine $M$ with a strict context bound $\kappa$, the maximum number of possible distinct, valid computational configurations is exactly bounded by $\mathcal{O}(|Q| \cdot \kappa \cdot |\Gamma|^\kappa)$.
\end{lemma}

\begin{proof}
A global configuration of a Turing machine $M$ at any discrete time step $t$ is fully specified by three independent parameters:
\begin{enumerate}
    \item \textbf{Internal State:} The current state $q \in Q$ of the finite control. There are exactly $|Q|$ distinct possibilities.
    \item \textbf{Tape Head Position:} The current physical position of the tape head. By Definition \ref{def:cbtm}, the machine immediately rejects if the head moves beyond $\kappa$ distinct cells. Thus, the head position $p$ is strictly confined to the interval $1 \le p \le \kappa$, yielding exactly $\kappa$ valid positions.
    \item \textbf{Tape Contents:} The sequence of symbols currently written on the active portion of the tape. There are at most $\kappa$ such cells initialized or subsequently modified. Each cell can independently hold any symbol drawn from the tape alphabet $\Gamma$. By elementary combinatorics, the number of distinct string permutations of length $\kappa$ over alphabet $\Gamma$ is $|\Gamma|^\kappa$.
\end{enumerate}

Since these three parameters constitute a complete representation of the machine's state, the total number of valid, non-rejecting configurations $\mathcal{N}$ is the Cartesian product of the individual parameter spaces:
\begin{equation}
\mathcal{N} = |Q| \times \kappa \times |\Gamma|^\kappa
\end{equation}
Because the state set $Q$, the context limit $\kappa$, and the tape alphabet $\Gamma$ are all strictly finite, non-negative integers, it follows that the entire configuration space is strictly bounded and finite, scaling exponentially only with respect to the context bound $\kappa$.
\end{proof}

The finiteness of the configuration space has profound implications for the decidability of execution within isolated Praxis agents.

\begin{theorem}[Decidability of CBTM Halting]
\label{thm:cbtm_halting}
The halting language for Context-Bounded Turing Machines, defined as 
$$HALT_{CBTM} = \{ \langle M, w \rangle \mid M \text{ is a CBTM that halts on input } w \}$$
resides strictly within the complexity class $\mathcal{P}_{SPACE}$ and is unreservedly decidable.
\end{theorem}

\begin{proof}
To prove decidability, we must construct a Turing machine that is guaranteed to halt and correctly decide whether a given CBTM $M$ halts on input $w$.
Let $\mathcal{N} = |Q| \times \kappa \times |\Gamma|^\kappa$ be the maximum absolute number of configurations of the CBTM $M$, as established in Lemma \ref{lem:finite_config}. 

We construct a universal deciding Turing machine $U$. The machine $U$ accepts as input the standard encoding $\langle M, w \rangle$. It initializes a simulation of $M$ executing on input $w$. 
Crucially, $U$ allocates a dedicated step counter $c$ on a secondary auxiliary tape, initialized to $0$. 

During each discrete simulated time step of $M$, $U$ performs the following operations:
\begin{enumerate}
    \item \textbf{Halting Check:} $U$ inspects the current simulated state of $M$. If $M$ has entered either the accept state $q_{accept}$ or the reject state $q_{reject}$, $U$ immediately halts and outputs \texttt{ACCEPT} or \texttt{REJECT} to match $M$'s behavior.
    \item \textbf{Context Violation Check:} $U$ computes the target destination of the tape head according to $M$'s transition function $\delta$. If the move places the head on the $(\kappa+1)$-th distinct cell, $U$ halts and outputs \texttt{REJECT}, simulating the intrinsic bounds of $M$.
    \item \textbf{Counter Increment:} $U$ increments the auxiliary step counter $c$ by exactly 1.
    \item \textbf{Infinite Loop Detection:} $U$ compares the current value of $c$ against the upper bound $\mathcal{N}$. If $c > \mathcal{N}$, the Pigeonhole Principle dictates that $M$ must have traversed at least one global configuration more than once. Because $M$ operates via a deterministic transition function, returning to a previously visited configuration implies that $M$ has entered an infinite, non-halting cycle. Recognizing this deterministic recurrence, $U$ intercepts the execution, halts, and outputs \texttt{REJECT}.
\end{enumerate}

Because the counter $c$ monotonically increases and the simulation is forcibly terminated if $c > \mathcal{N}$, the maximum number of simulation steps performed by $U$ is strictly bounded by $\mathcal{N}$. Therefore, $U$ is guaranteed to terminate on all possible inputs.

Regarding spatial complexity, the simulation tape for $M$ requires exactly $\mathcal{O}(\kappa)$ cells. The auxiliary counter tape requires enough cells to represent the integer $\mathcal{N}$ in binary, which is $\lceil \log_2 \mathcal{N} \rceil$. Expanding this:
\begin{equation}
\log_2(|Q| \cdot \kappa \cdot |\Gamma|^\kappa) = \log_2(|Q|) + \log_2(\kappa) + \kappa \log_2(|\Gamma|) \approx \mathcal{O}(\kappa \log |\Gamma|)
\end{equation}
The total space utilized by $U$ is thus bounded by a polynomial in $\kappa$. Consequently, the problem of deciding $HALT_{CBTM}$ is strictly decidable and intrinsically belongs to the complexity class $\mathcal{P}_{SPACE}$.
\end{proof}

\section{Decidability Algebra and the Receipt Verification Protocol}

In the Praxis architecture, as dictated by the stringent Rust core discipline standards, facts about the system's state or an agent's actions must be computed deterministically via cryptographic hashes. They must never be merely asserted. We formalize this mandate of cryptographic determinism through the framework of Decidability Algebra.

\begin{definition}[Receipt Function]
\label{def:receipt_function}
Let $\mathcal{A}$ represent the discrete topological space of all valid agentic actions, and let $\mathcal{S}$ represent the state space of the repository. A \textit{receipt function} $R: \mathcal{A} \times \mathcal{S} \rightarrow \{0, 1\}^{256}$ is an immutable mapping from a given state-action pair to a canonical BLAKE3 cryptographic hash. We rigorously define this function as:
\begin{equation}
R(a, s) = BLAKE3(\text{Serialize}(a) \oplus \text{Serialize}(s))
\end{equation}
where $\oplus$ represents concatenation, and $\text{Serialize}(x)$ is a deterministic function that transforms its operand into a strictly ordered N-Quads canonical byte stream, ensuring semantic graph isomorphisms map identically.
\end{definition}

By converting graph states into N-Quads and hashing them, Praxis offloads the NP-hard problem of graph isomorphism into an $\mathcal{O}(n)$ string comparison problem, protected by cryptographic collision bounds.

\begin{theorem}[Collision Resistance of Contextual Receipts]
\label{thm:collision_resistance}
Assuming the BLAKE3 hash function acts as an ideal pseudo-random oracle, the process of identifying two distinct pairs $(a_1, s_1)$ and $(a_2, s_2)$ such that $R(a_1, s_1) = R(a_2, s_2)$ requires an asymptotic computational time complexity bounded strictly below by $\Omega(2^{128})$.
\end{theorem}

\begin{proof}
This theorem evaluates the susceptibility of the receipt function to a generic Birthday Attack. Let the hash output space be uniform with a bit-length of $n=256$, giving a total output space of $N = 2^{256}$. 

The probability of finding at least one collision after generating $k$ independent hash receipts is derived from the generalized birthday problem:
\begin{equation}
P(\text{collision}) = 1 - \prod_{i=0}^{k-1} \left(1 - \frac{i}{N}\right)
\end{equation}
Using the standard Taylor series approximation $1 - x \approx e^{-x}$ for small $x$, we can bound the product:
\begin{equation}
P(\text{collision}) \approx 1 - e^{-\frac{k(k-1)}{2N}} \approx 1 - e^{-k^2 / (2 \cdot 2^n)}
\end{equation}
To achieve a collision probability of $P \ge 0.5$ (a successful forgery with better-than-even odds), an adversary must evaluate the receipt function $R$ approximately $k$ times, where:
\begin{equation}
0.5 = 1 - e^{-k^2 / (2 \cdot 2^{256})} \implies e^{-k^2 / 2^{257}} = 0.5 \implies \frac{k^2}{2^{257}} = \ln(2)
\end{equation}
\begin{equation}
k \approx \sqrt{2 \ln(2) \cdot 2^{256}} \approx 1.17 \times 2^{128} \implies \Omega(2^{128})
\end{equation}
Although the context bound $\kappa$ explicitly limits the maximum size of any single action $a$ or state $s$, the combinatoric domain $\mathcal{A} \times \mathcal{S}$ remains astronomically vast. Because $|\Gamma|^\kappa \gg 2^{256}$ for typical values of $\kappa$ (e.g., $100,000$ tokens over a vocabulary of $50,000$), the context bound does not artificially compress the preimage space. Consequently, algebraic inversion via exhaustive search remains strictly intractable, preserving the integrity of the receipt graph against manipulation.
\end{proof}

\subsection{Verification Complexity and N-Quads Ordering}

A receipt protocol is only useful if it can be rapidly verified. Praxis rejects verification protocols requiring superlinear time. 

\begin{lemma}[Linear Verification Bound]
\label{lem:linear_verification}
Verification of an N-Quads ordered receipt graph $G = (V, E)$ demands strictly $\mathcal{O}(|V| + |E|)$ computational time and $\mathcal{O}(1)$ auxiliary space during streaming.
\end{lemma}

\begin{proof}
By Praxis discipline, all topological facts forming the state $s$ must be pre-sorted into canonical N-Quads order before hashing. This canonicalization is guaranteed upstream.

A deterministic verifier $V_{chk}$ receives the serialized sequence of graph edges $S = \langle q_1, q_2, \dots, q_k \rangle$ iteratively via a streaming interface. For each contiguous string of bytes representing the fact $q_i$, $V_{chk}$ injects the bytes directly into the state machine of the BLAKE3 algorithm. 

BLAKE3 utilizes a Merkle tree internal structure and SIMD-optimized compression functions that process sequential blocks of data in $\mathcal{O}(L)$ time, where $L$ is the total byte length of the stream. Because the canonical N-Quads serialization maps the graph into a string of length $L \propto (|V| + |E|)$, the total time complexity required to compute the verifying hash is strictly bounded by $c \cdot (|V| + |E|)$, where $c$ is a small, hardware-dependent constant.

Crucially, because BLAKE3 computes rolling intermediate hashes, no prior facts $q_j$ (where $j < i$) need to be retained in auxiliary memory. The verifier maintains only the 256-bit internal chaining values of the hash state. Thus, the auxiliary space complexity is strictly bounded by $\mathcal{O}(1)$.
\end{proof}

\section{Algebraic Independence of Wall-Clock Time}

A critical invariant of the core team discipline states: \textit{"No wall clock in hash/receipt paths. Time only from graph OWL-Time literals."} To satisfy this invariant, we must mathematically prove that the receipt function $R$ acts as an endomorphism entirely independent of the physical time domain $\mathbb{T}$.

\begin{theorem}[Temporal Independence of Receipts]
\label{thm:temporal_independence}
Let $t_1, t_2 \in \mathbb{T}$ represent two distinct wall-clock times of execution. For any identically constructed state-action pair $(a, s)$, the execution trace guarantees that $R_{t_1}(a, s) = R_{t_2}(a, s)$.
\end{theorem}

\begin{proof}
Let the derivation of a receipt be formally represented by a deterministic function $\mathcal{F}: \mathcal{I} \rightarrow \mathcal{O}$, where $\mathcal{I}$ is the set of all systemic inputs to the hashing engine, and $\mathcal{O}$ is the resulting hash. 

By the strict architectural axioms of Praxis, the input space $\mathcal{I}$ is statically restricted. It encompasses only:
1. The serialized bytes of the action $a$.
2. The serialized bytes of the state $s$.
3. Static graph OWL-Time literals, which are immutable strings denoting logical temporal coordinates, rather than physical clock readings.

System time structures, such as \texttt{std::time::SystemTime} or \texttt{Instant::now()} in Rust, are syntactically and semantically banned from the execution paths computing $\mathcal{F}$. Consequently, $\mathcal{F}$ is a pure mathematical function with respect to the state transition; it possesses zero side effects and implicitly relies on zero external environmental variables.

Assume, for the sake of contradiction, that the receipts diverge across time: $R_{t_1}(a, s) \neq R_{t_2}(a, s)$. Because $\mathcal{F}$ is deterministic, divergent outputs necessitate divergent inputs. Thus, the total input sets at $t_1$ and $t_2$ must differ. 
However, the parameters $(a, s)$ and the embedded OWL-Time literals are stipulated to be identical. The only remaining explanation is the implicit inclusion of a hidden variable $h(t) \in \mathcal{I}$ that mutates as a function of wall-clock time $t$. 

But the Praxis invariant explicitly asserts the orthogonality of the input space and the time domain: $\mathcal{I} \cap \mathbb{T} = \emptyset$. The presence of $h(t) \in \mathcal{I}$ directly contradicts this axiom. 
Therefore, by contradiction, no such temporal divergence can exist. It rigorously follows that $R_{t_1}(a, s) = R_{t_2}(a, s)$ for all possible temporal evaluations $t_1, t_2$.
\end{proof}

\section{The Halting Problem in Unbounded Subagent Trees}

While isolated, singular agents operate within strictly bounded contexts and their halting properties are firmly decidable (as proven in Theorem \ref{thm:cbtm_halting}), an agent within the Praxis framework is granted the capability to recursively invoke subagents via system tools. We must analyze the theoretical limits of this composite, unbounded tree structure.

\begin{definition}[Subagent Tree Automaton (STA)]
A Subagent Tree Automaton (STA) is a nondeterministic abstract machine consisting of a network of individual CBTM nodes. An STA transitions not only by modifying its own local tape, but by executing a \texttt{spawn} instruction, which instantiates a new STA child node. The child is initialized with a tape derived from the parent's current state. The parent may suspend execution to wait for a child's message, or proceed in parallel. The global state of the computation is a dynamically expanding, asynchronous tree of active STAs.
\end{definition}

\begin{theorem}[Undecidability of STA Halting]
\label{thm:sta_halting}
The global halting problem for unbounded Subagent Tree Automata is undecidable, rendering the macroscopic Praxis system fully Turing-complete, despite the finite bounds of its microscopic components.
\end{theorem}

\begin{proof}
We proceed by reducing the halting problem of a standard Turing machine $M_{std}$ with an unbounded tape to the global halting problem of an STA. We must demonstrate that an STA can perfectly simulate the infinite tape of $M_{std}$.

Let $\kappa$ be the fixed context bound of a single STA node. We initialize a root STA node to simulate the control unit of $M_{std}$. Instead of using a single contiguous tape, we allow the tree structure to natively represent segments of the unbounded tape.
Each STA node represents a discrete, finite segment of the tape of exact length $\kappa$. 

The simulation algorithm proceeds as follows:
When the simulated head of $M_{std}$ attempts to cross the rightmost boundary of the current $\kappa$-length segment:
\begin{enumerate}
    \item \textbf{Allocation:} If the adjacent tape segment does not exist, the active STA issues an invocation call, spawning a new subagent to mathematically represent the newly allocated blank segment of tape.
    \item \textbf{Context Passing:} The parent STA transmits the current internal state $q$ of $M_{std}$ and the fact that the head is entering from the left edge via the subagent asynchronous message protocol.
    \item \textbf{Suspension:} The parent STA transitions to a suspended wait state, preserving its local tape segment but yielding active computation. Control is effectively delegated to the newly active child STA, which continues the simulation of $M_{std}$.
\end{enumerate}

Because the maximum topological depth of the subagent tree and the total number of instantiated nodes are strictly unbounded by the underlying system architecture, this recursive protocol successfully simulates an arbitrarily large, unbounded tape. 

If there existed an algorithm capable of deciding whether an arbitrary global STA tree eventually halts (i.e., whether all nodes eventually cascade into $q_{accept}$ or $q_{reject}$), we could trivially use that algorithm to decide $HALT_{TM}$ by wrapping any generic Turing machine in our STA simulation wrapper. Since $HALT_{TM}$ is famously proven to be undecidable by Turing, $HALT_{STA}$ must intrinsically be undecidable as well.
\end{proof}

This undecidability underscores the necessity of practical operational limits within the Praxis framework, such as maximum invocation depth constraints and temporal timeouts, to forcefully prune non-terminating subagent lineages.

\section{Algorithmic Surprises and O(n) Bounds}

A hallmark of failing software at scale is "hidden quadratic behavior," where innocent-looking abstractions cascade into $\mathcal{O}(n^2)$ bottlenecks. The core Praxis discipline categorically forbids this, requiring all algorithms operating over unbounded inputs to maintain strict linearity. We formally capture this requirement via the \textit{Linear Strictness Property}.

\begin{definition}[Linear Strictness]
An algorithm $A$, operating on an input of topological size $n$ and yielding an output of size $m$, is defined as \textit{linearly strict} if and only if:
\begin{enumerate}
    \item Its asymptotic worst-case time complexity satisfies $T(n) \le k \cdot n + c$.
    \item Its structural recursion depth satisfies $D(n) \le k' \cdot n + c'$.
\end{enumerate}
where $k, k', c, c'$ are strictly positive real constants.
\end{definition}

\begin{lemma}[Closure of Linear Strictness Under Bounded Composition]
\label{lem:linear_composition}
The sequential composition $A_2 \circ A_1$ of two linearly strict algorithms $A_1$ and $A_2$ remains linearly strict if and only if the output size function of $A_1$, denoted $m_1(n)$, scales no faster than $\mathcal{O}(n)$.
\end{lemma}

\begin{proof}
Let algorithm $A_1$ possess a time complexity $T_1(n) \le k_1 n + c_1$ and yield an output of size $m_1(n)$. 
Let algorithm $A_2$ possess a time complexity $T_2(m) \le k_2 m + c_2$.

The total time complexity of the sequentially composed function $A(x) = A_2(A_1(x))$ is simply the sum of their individual execution times:
\begin{equation}
T(n) = T_1(n) + T_2(m_1(n))
\end{equation}

\textbf{Forward Implication (Sufficiency):} Assume $m_1(n)$ is bounded by linearity, meaning there exist constants $c_3, c_4$ such that $m_1(n) \le c_3 n + c_4$. Substituting this bound into $T(n)$ yields:
\begin{equation}
T(n) \le (k_1 n + c_1) + k_2 (c_3 n + c_4) + c_2
\end{equation}
\begin{equation}
T(n) \le (k_1 + k_2 c_3) n + (c_1 + k_2 c_4 + c_2)
\end{equation}
By defining new constants $K = k_1 + k_2 c_3$ and $C = c_1 + k_2 c_4 + c_2$, we see that $T(n) \le K n + C$. This perfectly satisfies the definition of linear strictness.

\textbf{Reverse Implication (Necessity):} Conversely, assume $m_1(n)$ exhibits superlinear asymptotic growth, such as $m_1(n) = \Omega(n \log n)$ or $\Omega(n^2)$. Consequently, the time consumed by the second algorithm, $T_2(m_1(n))$, will similarly explode to $\Omega(n \log n)$ or $\Omega(n^2)$. The overall time $T(n)$ will therefore reflect this superlinear lower bound, inevitably violating the required $\mathcal{O}(n)$ constraint.

Thus, composition is safely linear if and only if intermediate payloads remain linearly bounded.
\end{proof}

This lemma is the mathematical bedrock justifying why intermediate memory representations (e.g., ASTs, Graph fragments) in the Praxis engine must never artificially duplicate nodes. Any representation mapping a linear input into a quadratic intermediate structure will critically fail during subsequent evaluations.

\section{Closed Vocabularies and the Refusal Manifold}

The Praxis engine operates strictly upon defined semantics. It mandates closed vocabularies characterized by specific prefixes (\texttt{wf:}, \texttt{hook:}, \texttt{prayer-kernel:}, \texttt{agent:}). Any token or predicate that falls outside this rigorously defined set must not trigger a panic or undefined behavior; instead, it must map to a highly structured \texttt{Refusal} variant. 

\begin{definition}[Refusal Manifold]
Let $\mathcal{V}_{known}$ denote the finite set of recognized vocabulary strings within the engine's lexical schema. For any arbitrary input predicate $p$, the evaluation function $E(p)$ maps the input into the disjoint union of outcomes: $Success \sqcup Refusal$. The specific subset of state-space evaluations that map cleanly to the $Refusal$ variant forms the topological construct known as the \textit{Refusal Manifold}, denoted $\mathcal{M}_R$.
\end{definition}

\begin{theorem}[Completeness of the Refusal Manifold]
\label{thm:refusal_completeness}
For any predicate string $p \notin \mathcal{V}_{known}$, the evaluation mapping guarantees that $E(p) \in \mathcal{M}_R$. Furthermore, this evaluation to a refusal completes in strictly $\mathcal{O}(|p|)$ time, precluding infinite traversal or panic states.
\end{theorem}

\begin{proof}
The evaluation engine parses the predicate $p$ sequentially as an array of characters. To guarantee $O(|p|)$ validation, we construct a deterministic finite automaton (DFA), specifically structurally implemented as a Trie (prefix tree). This DFA attempts to match the characters of $p$ against the legal elements of $\mathcal{V}_{known}$.

Because $\mathcal{V}_{known}$ is a statically closed and finite set, the maximum permissible depth of the DFA is bounded by a constant $L_{max} = \max_{v \in \mathcal{V}_{known}} |v|$.

The DFA transition logic enforces exactly one state transition per character of the input $p$. 
If, at any index up to $|p|$, the DFA attempts a transition on a character that does not exist in the current node's edges, it drops into a non-accepting trap state. This trap state signifies that the prefix of $p$ irreconcilably deviates from $\mathcal{V}_{known}$. Upon entering the trap state, the parser immediately halts and constructs a explicitly typed \texttt{Refusal} variant.

Because the parser inspects each character at most once, the maximum number of transitions is strictly $|p|$. The time complexity of routing the input into the Refusal Manifold is therefore strictly bounded by $\mathcal{O}(|p|)$. Furthermore, since Rust's static typing enforces that the return type of $E$ must be $Success$ or $Refusal$, and panics are entirely banned by the core team discipline, it is logically guaranteed that the output perfectly resides within $\mathcal{M}_R$.
\end{proof}

\section{Space Complexity of Error Paths vs Happy Paths}

A foundational, distinguishing axiom of the Praxis discipline states: \textit{"Error paths are tested as rigorously as happy paths"} and mandates that silent swallows (e.g., \texttt{.ok()}) are forbidden. Errors are first-class values. In this section, we analyze the space complexity divergence between successful standard execution (the "happy path") and refusal generation.

\begin{theorem}[Isomorphism of Space Complexity in Refusal Generation]
\label{thm:space_isomorphism}
Let $C_{happy}(n)$ denote the memory space complexity of a successfully evaluated Praxis rule operating on an input of size $n$. Let $C_{error}(n)$ denote the space complexity of the corresponding error path that terminates by generating a \texttt{Refusal} variant. It holds that either $C_{error}(n) = \Theta(C_{happy}(n))$ or $C_{error}(n) = \mathcal{O}(1)$.
\end{theorem}

\begin{proof}
Consider a generalized rule evaluation function $f: \Sigma^* \rightarrow \mathcal{R} \sqcup \mathcal{M}_R$, where $\mathcal{R}$ represents the successful state transformations and $\mathcal{M}_R$ is the Refusal Manifold. The execution profile dictates memory allocation. 

\textbf{Case 1 (Early Syntactic Rejection):} 
If the input $x \in \Sigma^*$ violates a shallow structural precondition—such as presenting an unknown vocabulary prefix parsed in Theorem \ref{thm:refusal_completeness}—the function $f(x)$ halts immediately. It terminates \textit{prior} to allocating the extensive internal memory required for deep semantic processing. 
In this scenario, $f(x)$ constructs the typed \texttt{Refusal} using a statically defined enum variant, requiring only a constant amount of stack space. Thus, $C_{error}(n) = \mathcal{O}(1)$, which trivially satisfies $\mathcal{O}(1) \le C_{happy}(n)$.

\textbf{Case 2 (Deep Semantic Rejection):} 
If the input $x$ is syntactically valid and passes initial lexing, but violates a deep logical invariant deep within the evaluation graph, the engine will have already parsed and allocated memory for the intermediate structures, such as the Abstract Syntax Tree (AST) or in-memory graph fragments. The size of these structures is fundamentally proportional to the input, demanding $\Theta(n)$ space.

Upon triggering the deep failure, the system must generate a \texttt{Refusal} variant. Because silent swallows are banned, the refusal must accurately capture the exact nature of the violation, often enclosing a subset of the currently allocated state or a pointer to the offending token. However, no panics occur; memory unwinding behaves predictably. The memory utilized at the moment of the refusal generation remains $\Theta(n)$ for the resident AST, plus an $\mathcal{O}(1)$ allocation for the refusal struct wrapper itself. Consequently, the peak memory footprint of the error path is exactly $C_{error}(n) = \Theta(n)$. 

Because the successful path $C_{happy}(n)$ for deep evaluations is similarly bounded below by $\Omega(n)$ and bounded above by $\mathcal{O}(n)$ (due to the mandates of linear strictness outlined in Lemma \ref{lem:linear_composition}), it inescapably follows that:
\begin{equation}
C_{error}(n) = \Theta(C_{happy}(n))
\end{equation}

This mathematical isomorphism proves that entering an error state does not induce uncontrolled memory bloat or memory leaks. The system maintains perfect structural symmetry in memory bounds regardless of whether the trajectory evaluates to success or rejection.
\end{proof}

\section{Conclusions on Intractability and Verifiability}

In this chapter, we have mapped the practical, highly optimized engineering tenets of the Praxis architecture to the rigorous framework of formal computational complexity. 

We established that while the halting of isolated, context-bounded agents is firmly decidable and operates within $\mathcal{P}_{SPACE}$, the recursive, unbounded invocation of the macroscopic subagent tree restores full Turing-completeness and, consequently, undecidability (Theorem \ref{thm:sta_halting}). 

Through Decidability Algebra, we bounded the critical verification step of the receipt protocol to linear $\mathcal{O}(n)$ time via canonical N-Quads sorting (Lemma \ref{lem:linear_verification}). This architectural choice explicitly averts the computationally disastrous $\mathcal{O}(n^2)$ subgraph-isomorphism checks that traditionally paralyze complex semantic network evaluations. 

Finally, by formalizing Linear Strictness and the Refusal Manifold, we proved that the engine scales predictably, and that error paths compute with the exact same deterministic efficiency and spatial symmetry as happy paths. These foundational mathematical limits are not mere theoretical exercises; they irrevocably inform, justify, and enforce the physical engineering constraints maintained by the core team standing gates.